%==============================================================================
% فصل اول: مقدمه، کلیات، بیان مسأله
%==============================================================================
\ChapterStart
\chapter{مقدمه، کلیات و بیان مسأله}
\label{ch:introduction}




\section{بیان مسأله}
\label{sec:problem-statement}


بیماری‌های قلبی-عروقی به عنوان علت اصلی مرگ‌ومیر در سطح جهان شناخته می‌شوند و بار اقتصادی، اجتماعی و بهداشتی بی‌سابقه‌ای را بر جوامع بشری و سیستم‌های سلامت تحمیل می‌کنند\cite{Palaniappan2026HeartStats}. بر اساس جدیدترین و جامع‌ترین گزارش آماری انجمن قلب آمریکا (\lr{AHA}) که در سال ۲۰۲۶ با عنوان "آمار بیماری‌های قلبی و سکته مغزی" منتشر گردیده است، بیماری‌های قلبی-عروقی مسبب بیش از $20.5$ میلیون مرگ سالانه در سراسر جهان هستند که این رقم نمایانگر یک‌سوم از کل مورتالیتی (مرگ‌ومیر) جهانی است\cite{Palaniappan2026HeartStats}. در ایالات متحده آمریکا، این آمار معادل وقوع یک مرگ ناشی از عوارض قلبی-عروقی در هر ۳۴ ثانیه برآورد شده است\cite{Palaniappan2026HeartStats}. در این میان، بیماری ایسکمیک قلبی\footnote{\lr{Ischemic heart disease}} و تظاهر حاد و تهدیدکننده حیات آن یعنی انفارکتوس حاد میوکارد\footnote{\lr{Acute myocardial infarction}} یا سکته قلبی، شایع‌ترین و مرگبارترین نوع این بیماری‌ها محسوب می‌گردد\cite{Palaniappan2026HeartStats,WHO2014NCD}. در ایران نیز، همگام با تغییرات اپیدمیولوژیک، گذار جمعیتی، شهرنشینی و تغییر سبک زندگی، انفارکتوس میوکارد به عنوان عامل اصلی مورتالیتی (مرگ‌ومیر) و موربیدیتی (ابتلا) شناخته می‌شود. داده‌های ثبت ملی انفارکتوس میوکارد ایران\footnote{\lr{Iranian Myocardial Infarction Registry}} (\lr{IMIR}) نشان می‌دهد که میزان بروز استانداردشده سنی سکته قلبی در ایران حدود $73.3$ در هر ۱۰۰ هزار نفر است و متأسفانه نرخ مرگ‌ومیر داخل بیمارستانی ناشی از آن رقمی در حدود $12.1$ درصد گزارش شده است\cite{Beyranvand2023IranMI}. سن ابتلا در جمعیت ایرانی به‌طور میانگین $61.2$ سال ارزیابی شده است که نشان‌دهنده ابتلای زودهنگام‌تر و درگیری سهم قابل‌توجهی از جمعیت فعال و نیروی کار جامعه در مقایسه با استانداردهای جهانی است\cite{Beyranvand2023IranMI}.

با وجود پیشرفت‌ در مداخلات فاز حاد مانند درمان‌های ترومبولیتیک، مداخله عروق کرونر از راه پوست (\lr{PCI}) و جراحی بای‌پس شریان کرونری (\lr{CABG}) که منجر به افزایش چشمگیر نرخ بقای بیماران در مرحله حاد سکته قلبی شده‌اند، این بیماران پس از ترخیص همچنان در معرض خطر بسیار بالای تجربه مجدد حوادث ایسکمیک، نارسایی قلبی و مرگ ناگهانی قلبی قرار دارند\cite{Byrne2023ESC}. به همین دلیل، اجرای راهبردهای پیشگیری ثانویه\footnote{\lr{Secondary prevention}}، به عنوان یک رکن جدایی‌ناپذیر و حیاتی در مدیریت بلندمدت بیماران پس از ترخیص شناخته می‌شود. پیشگیری ثانویه شامل مجموعه‌ای از مداخلات دارویی و غیردارویی است که با هدف کند کردن سیر پیشرفت آترواسکلروز، جلوگیری از ترومبوز مجدد و کاهش عوارض قلبی طراحی شده‌اند\cite{Byrne2023ESC,Virani2023ChronicCoronary}.

با این وجود، یکی از چالش‌برانگیزترین معضلات در طبابت روزمره و مدیریت بیماری‌های مزمن، شکاف عمیق میان دستورالعمل‌های استاندارد بالینی و رفتار واقعی بیماران است. اگرچه پایبندی به رژیم‌های دارویی پس از ترخیص تا حدودی در بیماران قلبی قابل‌قبول‌تر است، اما میزان پایبندی به اقدامات غیردارویی و اصلاحات سبک زندگی\footnote{\lr{Lifesyle modifications}}، که شامل رژیم غذایی سالم، فعالیت بدنی منظم، ترک دخانیات و مدیریت استرس است، به‌شدت نامطلوب و نگران‌کننده گزارش می‌شود\cite{Keenan2017Adherence,Kotseva2019EUROASPIRE}. مطالعات جهانی نشان داده‌اند که در ماه‌های پس از سکته قلبی، تنها ۲۹ درصد از بیماران موفق به دستیابی به اهداف بالینی مرتبط با هر سه فاکتور اصلی سبک زندگی (تغذیه، ورزش و وزن) می‌شوند؛ به‌طوری‌که ۶۶ درصد از آنان فاقد فعالیت بدنی کافی (کمتر از ۳۰ دقیقه در روز)، ۳۸ درصد مبتلا به چاقی مفرط، و ۱۹ درصد همچنان به مصرف دخانیات ادامه می‌دهند (که ۵۵ درصد از این افراد سیگاری‌های پیگیر و مقاوم به ترک هستند)\cite{Kotseva2019EUROASPIRE}. عدم تبعیت از این اصول پیشگیرانه پیامدهای ویرانگری دارد؛ مطالعات اثبات کرده‌اند که قطع زودهنگام و عدم پایبندی به اقدامات پیشگیری ثانویه می‌تواند نرخ بقای یک‌ساله بیماران را از $97.7$ درصد به $88.5$ درصد کاهش دهد\cite{Hammer2024Adherence}.

این عدم پایبندی یک پدیده تک‌عاملی نیست، بلکه تحت تأثیر شبکه‌ای پیچیده از عوامل فردی، اجتماعی، روانی و سیستماتیک قرار دارد\cite{Mohandas2025RecurrentMI}. در این راستا، عواملی نظیر سواد سلامت\footnote{\lr{Health literacy}} و درک از بیماری\footnote{\lr{Illness perception}} به عنوان دو فاکتور شناختی-رفتاری کلیدی در تعیین میزان تبعیت بیمار از پروتکل‌های پیشگیری ثانویه شناخته شده‌اند\cite{Foxwell2013IllnessPerceptions}. بیماران مراجعه‌کننده به کلینیک قلب و عروق شهید مدنی خرم‌آباد، به عنوان یکی از مراکز ارجاعی مهم در غرب کشور، به واسطه ساختار فرهنگی، اجتماعی و اقتصادی خاص منطقه، ممکن است الگوهای متفاوتی از درک بیماری، سواد سلامت و دسترسی به منابع توانبخشی را تجربه کنند. فقدان داده‌های بومی و مستند در خصوص میزان دقیق پایبندی بیماران این منطقه به اقدامات غیردارویی و عدم شناخت کافی از موانع شناختی آنان، منجر به کاهش اثربخشی برنامه‌های توانبخشی و افزایش نرخ بستری‌های مجدد می‌گردد. از این رو، انجام این مطالعه با هدف بررسی میزان پایبندی به اقدامات پیشگیری ثانویه غیردارویی و تبیین نقش عوامل مرتبط نظیر سواد سلامت و درک از بیماری در بیماران مبتلا به انفارکتوس میوکارد در خرم‌آباد، یک ضرورت بالینی و بهداشتی انکارناپذیر است تا زمینه برای طراحی مداخلات آموزشی و حمایتی شخصی‌سازی‌شده و مبتنی بر شواهد در این منطقه فراهم گردد.

\section{کلیات و چارچوب نظری}
\subsection{ اپیدمیولوژی و کلیات سکته قلبی و پیشگیری ثانویه}


انفارکتوس حاد میوکارد نتیجه انسداد ناگهانی شریان‌های کرونری و قطع جریان خون به بخشی از عضله قلب است که منجر به ایسکمی و نکروز بافتی می‌شود\cite{Byrne2023ESC,Virani2023ChronicCoronary}. از منظر اپیدمیولوژیک، بروز سکته قلبی در سطح جهانی و ملی همچنان به عنوان یک دغدغه اصلی سلامت عمومی مطرح است. گزارش سال ۲۰۲۶ انجمن قلب آمریکا (\lr{AHA}) نشان می‌دهد که شیوع عوامل خطر اصلی بیماری‌های قلبی-عروقی شامل پرفشاری خون، چاقی مفرط، و دیابت به سطوح بی‌سابقه‌ای در سطح جهانی رسیده است\cite{Palaniappan2026HeartStats}. آمارهای ملی ثبت سکته قلبی (\lr{IMIR}) نشان می‌دهند که فاکتورهای خطر کلاسیک در میان بیماران ایرانی مبتلا به سکته قلبی با فرکانس بالایی حضور دارند. شیوع پرفشاری خون در این بیماران تا 50 درصد (در برخی گزارش‌ها 44 درصد)، مصرف مستمر دخانیات $26.2$ درصد، دیابت نوع دو $22.2$ درصد و دیس‌لیپیدمی $17.8$ درصد گزارش شده است\cite{Beyranvand2023IranMI}. این پروفایل پرخطر متابولیک، فیزیولوژی عروق را مستعد تشکیل مجدد پلاک‌ها و وقوع حوادث مکرر قلبی می‌نماید. در حالی که در کشورهای توسعه‌یافته نرخ بروز و مرگ‌ومیر قلبی به دلیل اجرای برنامه‌های جامع پیشگیری اولیه روندی نزولی یا ثابت به خود گرفته است، در کشورهای با درآمد متوسط و پایین، از جمله ایران، این روند همچنان صعودی است\cite{WHO2014NCD,Beyranvand2023IranMI}.

پیشگیری ثانویه پس از انفارکتوس میوکارد به مجموعه‌ای از اقدامات هماهنگ و چندوجهی اطلاق می‌شود که برای بازگرداندن بیمار به بالاترین سطح ممکن از سلامت جسمی، روانی و اجتماعی و جلوگیری از پیشرفت بیماری آترواسکلروتیک طراحی شده است. دستورالعمل‌های معتبر بین‌المللی نظیر کالج کاردیولوژی آمریکا\footnote{\lr{American College of Cardiology}} (\lr{ACC})، انجمن قلب آمریکا\footnote{\lr{American Heart Association}} (\lr{AHA}) و انجمن قلب اروپا\footnote{\lr{European Society of Cardiology}} (\lr{ESC}) تأکید ویژه‌ای بر ترکیب درمان‌های دارویی (شامل درمان ضدپلاکت دوگانه، استاتین‌ها با دوز بالا، بتا-بلاکرها و مهارکننده‌های سیستم رنین-آنژیوتانسین) با اصلاحات جدی و پایدار در سبک زندگی دارند\cite{Byrne2023ESC,Virani2023ChronicCoronary}. رویکرد نوین گایدلاین‌ها تأکید دارد که درمان دارویی به‌تنهایی برای مهار فرآیند پاتولوژیک آترواسکلروز کافی نیست و مداخلات غیردارویی نه تنها اثربخشی داروها را تقویت می‌کنند، بلکه خود به صورت مستقل اثرات ضدالتهابی، آنتی‌ترومبوتیک و تثبیت‌کننده پلاک‌های عروقی دارند\cite{Byrne2023ESC,Keenan2017Adherence}.


\subsection{اقدامات غیردارویی در پیشگیری ثانویه}


مداخلات غیردارویی در بیماران پس از سکته قلبی به عنوان بنیان توانبخشی قلبی (\lr{Cardiac Rehabilitation}) شناخته می‌شوند که بر سه محور اساسی رژیم غذایی، فعالیت بدنی و ترک دخانیات استوارند:

\textbf{رژیم غذایی و تغذیه:} تغییر در عادات تغذیه‌ای و حفظ آن در طولانی‌مدت، یکی از چالش‌برانگیزترین اجزای پیشگیری ثانویه است. گایدلاین‌های \lr{ACC/AHA} تبعیت از الگوهای غذایی دوستدار قلب، به‌ویژه رژیم غذایی مدیترانه‌ای یا رژیم \lr{DASH}\footnote{\lr{Dietary Approaches to Stop Hypertention}} را توصیه می‌کنند\cite{Virani2023ChronicCoronary,Estruch2018Mediterranean}. این رژیم‌ها بر مصرف بالای سبزیجات، میوه‌های تازه، غلات کامل، حبوبات، پروتئین‌های بدون چربی (مانند ماهی که منبع غنی اسیدهای چرب امگا-۳ است) و روغن‌های گیاهی غیراشباع تأکید دارند، در حالی که مصرف چربی‌های ترانس، اسیدهای چرب اشباع، کربوهیدرات‌های تصفیه‌شده، گوشت‌های فرآوری‌شده و سدیم باید به حداقل برسد\cite{Estruch2018Mediterranean}. مطالعات نشان می‌دهد که رعایت چنین رژیم‌هایی از طریق بهبود پروفایل لیپیدی، کاهش استرس اکسیداتیو، سرکوب مسیرهای التهابی سیستمیک (نظیر کاهش سطوح \lr{hs-CRP}) و ارتقای آزادسازی نیتریک اکساید (\lr{NO}) اندوتلیال، خطر حوادث ایسکمیک مکرر را به شکل معناداری کاهش می‌دهد\cite{Estruch2018Mediterranean}.

\textbf{فعالیت بدنی و بازتوانی قلبی: } استراحت مطلق و بی‌تحرکی پس از سکته قلبی، پارادایمی منسوخ شده است. دستورالعمل‌های فعلی توصیه می‌کنند که بیماران ترخیص‌شده باید به سرعت تحت برنامه‌های ساختاریافته بازتوانی قلبی قرار گیرند\cite{Virani2023ChronicCoronary}. انجام حداقل ۱۵۰ دقیقه فعالیت هوازی با شدت متوسط در هفته (یا ۷۵ دقیقه فعالیت شدید) ، به همراه تمرینات مقاومتی سبک، به عنوان هدف استاندارد تعریف شده است\cite{Virani2023ChronicCoronary}. فعالیت بدنی منظم از طریق تحریک آنژیوژنز (رگ‌زایی)، بهبود پاسخ عصب واگ\footnote{\lr{Vagal tone}}، ایجاد گردش خون جانبی\footnote{\lr{Collateral circulation}} و کاهش مقاومت به انسولین، مورتالیتی قلبی-عروقی را تا حدود ۲۵ الی ۳۰ درصد کاهش می‌دهد\cite{Virani2023ChronicCoronary,Kotseva2019EUROASPIRE}.

\textbf{ترک دخانیات:} قطع مصرف تنباکو به عنوان مقرون‌به‌صرفه‌ترین و مؤثرترین مداخله منفرد در پیشگیری ثانویه شناخته می‌شود. استمرار مصرف سیگار پس از انفارکتوس میوکارد باعث تحریک مداوم سیستم سمپاتیک، تشدید اسپاسم شریان‌های کرونر، کاهش ظرفیت حمل اکسیژن توسط هموگلوبین (به واسطه کربوکسی‌هموگلوبین)، افزایش ویسکوزیته خون و پیشرفت سریع‌تر آترواسکلروز می‌شود\cite{Gao2025SmokingCessation}.  مطالعات هم‌گروهی اثبات کرده‌اند که ترک کامل سیگار پس از \lr{AMI} می‌تواند خطر مرگ‌ومیر قلبی را تا ۳۶ درصدکاهش دهد\cite{Gao2025SmokingCessation}.

\textbf{مدیریت استرس و سلامت روان\footnote{\lr{Psychosocial Stress Management}}:} افسردگی، اضطراب مزمن و استرس‌های روانی-اجتماعی در فاز پس از سکته قلبی بسیار شایع بوده و به عنوان عوامل خطر مستقل برای پیش‌آگهی نامطلوب شناخته می‌شوند. این عوامل از طریق فعال‌سازی محور هیپوتالاموس-هیپوفیز-آدرنال (\lr{HPA})، افزایش ترشح کاتکول‌آمین‌ها و کورتیزول، و افزایش سطح سایتوکین‌های ‌التهابی (نظیر \lr{IL-6} و \lr{hs-CRP}) ، سیستم قلبی-عروقی را تحت فشار قرار می‌دهند\cite{Vaccarino2020Depression}. لذا غربالگری دوره‌ای سلامت روان، ارجاع به مشاوره‌های روان‌شناختی و آموزش تکنیک‌های مدیریت استرس، جزء ضروری در پروتکل‌های جامع پیشگیری ثانویه محسوب می‌گردد\cite{Vaccarino2020Depression}.

\subsection{مفهوم پایبندی به درمان و اقدامات خودمراقبتی}

در متون آکادمیک پزشکی، مفهوم پایبندی\footnote{\lr{Adherence}} تکامل‌یافته‌ی واژه قدیمی‌تر تمکین\footnote{\lr{Compliance}} است. در حالی که تمکین به معنای اطاعت منفعلانه بیمار از دستورات پزشک بود، سازمان جهانی بهداشت (\lr{WHO}) پایبندی را به عنوان "میزانی که رفتار یک فرد (از جمله مصرف دارو، پیروی از رژیم غذایی و اجرای تغییرات سبک زندگی) با توصیه‌های توافق‌شده بین بیمار و ارائه‌دهنده خدمات بهداشتی مطابقت دارد" تعریف می‌کند\cite{Chaudri2004AdherenceWHO}. این تغییر پارادایم، بر نقش فعال و تصمیم‌گیرنده بیمار در مدیریت بیماری خویش (خودمراقبتی) تأکید دارد.

در پژوهش‌های بالینی، سنجش میزان پایبندی به مداخلات غیردارویی به مراتب پیچیده‌تر و چالش‌برانگیزتر از پایش مصرف داروها ارزیابی می‌شود. با این حال، ارزیابی پایبندی به سبک زندگی اغلب از طریق پرسشنامه‌های خودگزارش‌دهی استانداردشده، مصاحبه‌های بالینی و یا بررسی دستیابی بیمار به اهداف فیزیولوژیک (نظیر کنترل فشار خون زیر ۹۰/۱۴۰ میلی‌متر جیوه، کنترل \lr{BMI} در محدوده $18.5$ تا $24.9$ کیلوگرم بر مترمربع، تثبیت پروفایل لیپیدی و سطح \lr{HbA1c} در بیماران دیابتی) انجام می‌گیرد\cite{Kotseva2019EUROASPIRE,Chaudri2004AdherenceWHO}. شواهد اپیدمیولوژیک نشان می‌دهد که در حالی که نرخ پایبندی دارویی در شش ماه نخست پس از ترخیص ممکن است به بالای ۹۰ درصد برسد، نرخ استمرار پایبندی به توصیه‌های رژیم غذایی مناسب و ورزش منظم در بسیاری از جوامع، حتی با وجود آموزش‌های اولیه، به ندرت از ۳۰ تا ۵۰ درصد فراتر می‌رود\cite{He2026CardiacRehab}.

\subsection{عوامل مرتبط با پایبندی: سواد سلامت}

سواد سلامت\footnote{\lr{Health Literacy}} یکی از مهم‌ترین پیش‌بینی‌کننده‌های مستقل رفتارهای بهداشتی و پیامدهای درمانی در بیماران مبتلا به بیماری‌های مزمن است. سواد سلامت فراتر از توانایی ساده خواندن و نوشتن است و به عنوان "ظرفیت شناختی و مهارت‌های اجتماعی که انگیزه و توانایی افراد را برای دسترسی، درک پردازش و استفاده از اطلاعات بهداشتی به منظور ارتقا و حفظ سلامت تعیین می‌کند"، تعریف می‌شود\cite{Sorensen2015HealthLiteracy}.

در بستر پیشگیری ثانویه از سکته قلبی، بیمار با حجم انبوهی از اطلاعات پیچیده پزشکی، دستورالعمل‌های دارویی و ضرورت تغییرات رادیکال در سبک زندگی مواجه می‌شود. سواد سلامت در سه سطح عملکردی (توانایی خواندن برچسب‌های دارویی و غذایی)، تعاملی (توانایی برقراری ارتباط مؤثر با تیم درمان برای پرسش و رفع ابهام) و انتقادی (توانایی تحلیل اطلاعات بهداشتی و تصمیم‌گیری آگاهانه) نقش‌آفرینی می‌کند\cite{Sorensen2015HealthLiteracy,Zakeri2023HealthLiteracy}. بیمارانی که از سواد سلامت پایینی برخوردارند، در درک مکانیسم بیماری خود و منطق نهفته در توصیه‌های پیشگیرانه دچار مشکل می‌شوند. برای مثال، چنین بیمارانی ممکن است نتوانند ارتباط مستقیم میان مصرف بالای سدیم و تشدید نارسایی قلبی را درک کنند، یا مفهوم "شدت متوسط" در ورزش هوازی را به درستی تفسیر نمایند. مطالعات متعدد ثابت کرده‌اند که سواد سلامت پایین به‌طور مستقیم با عدم تبعیت از رژیم غذایی، ناتوانی در مدیریت دوز داروها، کاهش مشارکت در برنامه‌های بازتوانی قلبی و افزایش مکرر نرخ بستری‌های اورژانسی مرتبط است\cite{He2026CardiacRehab,Zakeri2023HealthLiteracy}. در واقع، سطح تحصیلات و درآمد ماهیانه در کنار ظرفیت سواد سلامت، تعیین‌کننده‌های بنیادین پایبندی بیماران به برنامه‌های بازتوانی محسوب می‌گردند\cite{He2026CardiacRehab}.

\subsection{عوامل مرتبط با پایبندی: درک از بیماری}

علاوه بر توانایی‌های شناختی (سواد سلامت)، باورها و برداشت‌های روانی بیمار از بیماری خود، که تحت عنوان «درک از بیماری»\footnote{\lr{Illness Perception}} شناخته می‌شود، نقش تعیین‌کننده‌ای در شکل‌گیری رفتارهای خودمراقبتی و پایبندی طولانی‌مدت دارد. مفهوم درک از بیماری بر پایه «مدل خودتنظیمی عقل سلیم لِوِنتال»\footnote{\lr{Leventhal's Common-Sense Model of Self-Regulation}} بنا شده است\cite{Leventhal2016CSM}. بر اساس این مدل روان‌شناختی، هنگامی که فرد با یک تهدید سلامتی نظیر سکته قلبی مواجه می‌شود، برای تفسیر این رویداد و سازگاری با آن، یک بازنمایی ذهنی\footnote{\lr{Mental Representation}} از بیماری در ذهن خود می‌سازد\cite{Leventhal2016CSM,Broadbent2015BriefIPQ}. این بازنمایی دارای پنج بعد شناختی اساسی است:

\begin{enumerate}

\item شناسه (\lr{Identity}): باور بیمار در مورد نام بیماری و علائم مرتبط با آن.

\item علت (\lr{Cause}): برداشت شخصی بیمار از عوامل ایجادکننده سکته قلبی (مانند استرس، ژنتیک، یا شانس).

\item سیر زمانی (\lr{Timeline}): درک بیمار از مدت‌زمان ماندگاری بیماری (حاد در برابر مزمن).

\item پیامدها (\lr{Consequences}): ارزیابی بیمار از تأثیرات و عوارض بیماری بر ابعاد مختلف زندگی شخصی، اقتصادی و خانوادگی.

\item کنترل و درمان‌پذیری (\lr{Control/Cure}): میزان باور بیمار به توانمندی شخصی خود (کنترل درونی) و اثربخشی درمان‌های پزشکی (کنترل بیرونی) در مهار بیماری\cite{Leventhal2016CSM,French2006CardiacRehab}.

\end{enumerate}

یکی از بزرگ‌ترین موانع در مسیر پایبندی به پیشگیری ثانویه، برداشت‌های غلط در بعد «سیر زمانی» و «درمان‌پذیری» است. شواهد علمی نشان می‌دهند که نرخ عدم پایبندی به داروها و اصلاحات سبک زندگی در بیماران قلبی بین ۳۳ تا بیش از ۵۵ درصد در نوسان است و این مسئله قویا با ادراکات بیمار پیوند دارد\cite{Mobini2023IllnessPerception}. بسیاری از بیماران پس از انجام موفقیت‌آمیز استنت‌گذاری (\lr{PCI})، سکته قلبی را یک رویداد "حاد" تلقی می‌کنند که اکنون به طور کامل "درمان" و برطرف شده است\cite{French2006CardiacRehab}. این درک نادرست که با ماهیت مزمن و پیشرونده آترواسکلروز در تضاد است، انگیزه بیمار را برای پذیرش و استمرار تغییرات دشوار غیردارویی در سبک زندگی (نظیر پرهیز غذایی دائمی یا ترک سیگار) به شدت سرکوب می‌کند. در مقابل، بیمارانی که درک صحیحی از مزمن بودن بیماری خود دارند و همزمان از حس بالای «کنترل‌پذیری» برخوردارند، بیماری را پدیده‌ای قابل‌مدیریت از طریق رفتار فعالانه خویش می‌پندارند. مطالعات روان‌تنی\footnote{\lr{Psychosomatic}} به وضوح نشان داده‌اند که درک مثبت از قابلیت کنترل بیماری و حس بالای خودکارآمدی\footnote{\lr{Self efficacy}}، قویا با سطح بالاتری از پایبندی به تغییرات سبک زندگی و مشارکت مستمر در برنامه‌های بازتوانی قلب همبستگی دارند\cite{Mobini2023IllnessPerception,Aujla2016IllnessBeliefs}. لذا، شناسایی این الگوهای ادراکی در بیماران ترخیص‌شده از اهمیت بالایی برخوردار است تا بتوان با مداخلات شناختی-رفتاری، باورهای مخرب را پیش از تثبیت، اصلاح نمود\cite{Mobini2023IllnessPerception,Aujla2016IllnessBeliefs}.


\section{اهداف پژوهش}
\label{sec:objectives}

\subsection{هدف کلی}


تعیین میزان پایبندی به اقدامات پیشگیری ثانویه غیردارویی و عوامل مرتبط با آن در بیماران سکته قلبی مراجعه کننده به درمانگاه قلب شهید مدنی شهر خرم آباد در نیمسال دوم 1404


\subsection{اهداف اختصاصی}
\begin{enumerate}[label=\arabic*.]
    
\item تعیین میزان پایبندی به ترک دخانیات در بیماران مورد مطالعه

\item تعیین میزان پایبندی به رژیم غذایی سالم (بر اساس معیار \lr{MEDAS}) در بیماران مورد مطالعه

\item تعیین میزان پایبندی به فعالیت بدنی منظم (بر اساس معیار \lr{IPAQ}) در بیماران مورد مطالعه

\item تعیین وضعیت مدیریت استرس و افسردگی (بر اساس معیار \lr{PHQ-4}) در بیماران مورد مطالعه

\item تعیین میزان پایبندی کلی به اقدامات پیشگیری ثانویه غیردارویی در بیماران مورد مطالعه

\item تعیین ارتباط سن با میزان پایبندی به اقدامات پیشگیری ثانویه غیردارویی

\item تعیین ارتباط جنسیت با میزان پایبندی به اقدامات پیشگیری ثانویه غیردارویی

\item تعیین ارتباط سطح تحصیلات با میزان پایبندی به اقدامات پیشگیری ثانویه غیردارویی

\item تعیین ارتباط وضعیت تأهل با میزان پایبندی به اقدامات پیشگیری ثانویه غیردارویی

\item تعیین ارتباط محل سکونت با میزان پایبندی به اقدامات پیشگیری ثانویه غیردارویی

\item تعیین ارتباط سطح سواد سلامت با میزان پایبندی به اقدامات پیشگیری ثانویه غیردارویی

\item تعیین ارتباط درک از بیماری (\lr{Illness Perception}) با میزان پایبندی به اقدامات پیشگیری ثانویه غیردارویی

\item تعیین ارتباط زمان سپری‌شده از سکته قلبی با میزان پایبندی به اقدامات پیشگیری ثانویه غیردارویی

\item تعیین ارتباط اقدام درمانی انجام شده (\lr{PCI/CABG/Medical}) با میزان پایبندی به اقدامات پیشگیری ثانویه غیردارویی

\item تعیین ارتباط تعداد بیماری‌های همراه (\lr{Comorbidities}) با میزان پایبندی به اقدامات پیشگیری ثانویه غیردارویی

\item تعیین ارتباط سابقه خانوادگی ابتلا به بیماری قلبی عروقی با میزان پایبندی به اقدامات پیشگیری ثانویه غیردارویی

\item تعیین ارتباط وضعیت اقتصادی (ادراک شده) با میزان پایبندی به اقدامات پیشگیری ثانویه غیردارویی

\item تعیین ارتباط وضعیت شغلی با میزان پایبندی به اقدامات پیشگیری ثانویه غیردارویی

\item تعیین ارتباط وضعیت پوشش بیمه با میزان پایبندی به اقدامات پیشگیری ثانویه غیردارویی

\item تعیین ارتباط پایبندی به درمان دارویی (بر اساس \lr{MMAS-4} ) با میزان پایبندی به اقدامات پیشگیری ثانویه غیردارویی

\item تعیین ارتباط شرکت در بازتوانی قلبی با میزان پایبندی به اقدامات پیشگیری ثانویه غیردارویی

\end{enumerate}



\section{پرسش‌های پژوهش}
\label{sec:research-questions}
\begin{enumerate}[label=\arabic*.]

\item میزان پایبندی به ترک دخانیات در بیماران مورد مطالعه چقدر است؟

\item میزان پایبندی به رژیم غذایی سالم (بر اساس معیار \lr{MEDAS}) در بیماران مورد مطالعه چقدر است؟

\item میزان پایبندی به فعالیت بدنی منظم (بر اساس معیار \lr{IPAQ}) در بیماران مورد مطالعه چقدر است؟

\item وضعیت مدیریت استرس و افسردگی (بر اساس معیار \lr{PHQ-4}) در بیماران مورد مطالعه چگونه است؟

\item میزان پایبندی کلی به مجموعه اقدامات پیشگیری ثانویه غیردارویی (شامل هر ۴ رکن فوق) در بیماران مورد مطالعه چقدر است؟

\end{enumerate}

\section{فرضیات پژوهش}
\label{sec:hypotheses}
\begin{enumerate}[label=\arabic*.]

\item سن با میزان پایبندی به اقدامات پیشگیری ثانویه غیردارویی ارتباط دارد.

\item جنسیت با میزان پایبندی به اقدامات پیشگیری ثانویه غیردارویی ارتباط دارد.

\item سطح تحصیلات با میزان پایبندی به اقدامات پیشگیری ثانویه غیردارویی ارتباط دارد.

\item وضعیت تأهل با میزان پایبندی به اقدامات پیشگیری ثانویه غیردارویی ارتباط دارد.

\item محل سکونت (شهری/روستایی) با میزان پایبندی به اقدامات پیشگیری ثانویه غیردارویی ارتباط دارد.

\item سطح سواد سلامت با میزان پایبندی به اقدامات پیشگیری ثانویه غیردارویی ارتباط دارد.

\item نمره درک از بیماری (\lr{Illness Perception}) با میزان پایبندی به اقدامات پیشگیری ثانویه غیردارویی ارتباط دارد.

\item زمان سپری‌شده از سکته قلبی با میزان پایبندی به اقدامات پیشگیری ثانویه غیردارویی ارتباط دارد.

\item نوع اقدام درمانی انجام شده (\lr{PCI}, \lr{CABG}, \lr{Medical}) با میزان پایبندی به اقدامات پیشگیری ثانویه غیردارویی ارتباط دارد.

\item تعداد بیماری‌های همراه (\lr{Comorbidities}) با میزان پایبندی به اقدامات پیشگیری ثانویه غیردارویی ارتباط دارد.

\item سابقه خانوادگی ابتلا به بیماری قلبی-عروقی با میزان پایبندی به اقدامات پیشگیری ثانویه غیردارویی ارتباط دارد.

\item وضعیت اقتصادی (ادراک شده) با میزان پایبندی به اقدامات پیشگیری ثانویه غیردارویی ارتباط دارد.

\item وضعیت شغلی با میزان پایبندی به اقدامات پیشگیری ثانویه غیردارویی ارتباط دارد.

\item وضعیت پوشش بیمه با میزان پایبندی به اقدامات پیشگیری ثانویه غیردارویی ارتباط دارد.

\item پایبندی به درمان دارویی (بر اساس نمره \lr{MMAS-4}) با میزان پایبندی به اقدامات پیشگیری ثانویه غیردارویی ارتباط دارد.

\item سابقه شرکت در بازتوانی قلبی با میزان پایبندی به اقدامات پیشگیری ثانویه غیردارویی ارتباط دارد.

\end{enumerate}

